\section{Limitations and Drawbacks of Our Work}
\label{sec:limitations}

While the PrimeFusion FANET framework offers significant advantages in terms of security and efficiency, it is essential to acknowledge its inherent limitations and the trade-offs involved in its design. A balanced scientific assessment requires a clear articulation of the system's weaknesses, which provides context for its applicability and guides future research. This section explicitly discusses the primary drawbacks of our proposed approach.

\subsection{Computational and Energy Overhead}
Despite the design of a lightweight consensus mechanism, the integration of blockchain technology inevitably introduces computational and energy overhead compared to non-blockchain-based systems. The processes of cryptographic signing, transaction verification, and ledger maintenance, although optimized in PoC, still consume CPU cycles and energy resources. For small, energy-constrained UAVs with limited processing power, this overhead could potentially impact flight duration and the capacity to perform other mission-critical computations. The CBOR compression layer, while beneficial for bandwidth, also adds a small computational burden for encoding and decoding data. Therefore, a careful cost-benefit analysis is required for any given mission to determine if the enhanced security justifies the additional resource consumption.

\subsection{Latency in Consensus and Transaction Finality}
The Proof-of-Cooperation consensus, while significantly faster than PoW, is not instantaneous. There is an inherent latency associated with block creation, propagation, and validation across the network. For highly dynamic control loops or applications requiring near-real-time data synchronization (e.g., coordinated acrobatic maneuvers), even the sub-50ms latency achieved by PrimeFusion might be prohibitive. The concept of \textit{transaction finality} in blockchain means that a transaction is only considered irreversible after it has been included in a block and confirmed by a sufficient number of subsequent blocks. This delay, while crucial for security, represents a fundamental trade-off against the immediacy of data availability that is often desired in tactical communication systems.

\subsection{Scalability Constraints}
Our framework was successfully validated with up to 30 UAVs, demonstrating good scalability and maintaining a high Packet Delivery Ratio (PDR). However, the performance of any blockchain-based system will eventually degrade as the number of nodes and the transaction volume increase. In ultra-dense swarm scenarios involving hundreds or thousands of UAVs, the overhead of maintaining a single, unified blockchain could become untenable. The communication burden of broadcasting every transaction to every node would lead to network congestion and increased consensus latency. While techniques such as sharding or hierarchical blockchain architectures could potentially address these issues, they are not part of the current PrimeFusion framework and would introduce significant additional complexity.

\subsection{Dependence on Connectivity}
Like all distributed consensus systems, the PoC mechanism is predicated on the existence of a reasonably well-connected network. In scenarios where the FANET becomes partitioned into several disconnected sub-networks due to distance, obstacles, or jamming, each partition might start to build its own version of the blockchain (a fork). While the framework is designed to resolve these forks once connectivity is restored, frequent or prolonged network partitions could undermine the integrity of the global ledger and lead to inconsistencies in the mission data. This reliance on connectivity makes the system vulnerable in contested or communication-degraded environments.

